\chapter{Ma Trận Vuông}
\hypertarget{localtoc\thechapter}{}
\localtoc

\begin{problem}
    Viết hàm nhập, xuất ma trận vuông các số nguyên

Hàm duyệt các phần tử trên đường chéo chính

Hàm duyệt các phần tử thuộc tam giác trên đường chéo chính

Hàm duyệt các phần tử thuộc tam giác dưới đường chéo chính

Hàm duyệt các phần tử trên đường chéo phụ

Hàm duyệt các phần tử thuộc tam giác trên đường chéo phụ

Hàm duyệt các phần tử thuộc tam giác dưới đường chéo phụ
\end{problem}
\begin{problem}
    Viết hàm nhập, xuất ma trận vuông các số thực
\end{problem}
\begin{problem}
    Viết hàm tìm giá trị lớn nhất trong ma trận vuông các số thực
\end{problem}
\begin{problem}
    Viết hàm kiểm tra trong ma trận vuông các số nguyên có tồn tại giá trị chẵn nhỏ hơn 2015 hay không
\end{problem}
\begin{problem}
    Viết hàm đếm số lượng số nguyên tố trong ma trận vuông các số nguyên
\end{problem}
\begin{problem}
    Viết hàm tính tổng các giá trị âm trong ma trận vuông các số thực
\end{problem}
\begin{problem}
    Viết hàm sắp xếp ma trận vuông các số thực tăng dần từ trên xuống dưới và từ trái sang phải
\end{problem}


\section{Kĩ thuật tính toán}
\begin{problem}
    Tổng phần tử thuộc ma trận tam giác trên (không tính đường chéo) trong ma trận vuông
\end{problem}
\begin{problem}
    Tổng phần tử thuộc ma trận tam giác dưới (không tính đường chéo) trong ma trận vuông
\end{problem}
\begin{problem}
    Tổng phần tử trên đường chéo chính
\end{problem}
\begin{problem}
    Tổng phần tử trên đường chéo phụ
\end{problem}
\begin{problem}
    Tổng phần tử chẵn nằm trên biên
\end{problem}

\section{Kỹ thuật đặt lính canh}
\begin{problem}
    Tìm max trong ma trận tam giác trên
\end{problem}
\begin{problem}
    Tìm min trong ma trận tam giác dưới
\end{problem}
\begin{problem}
    Tìm max trên đường chéo chính
\end{problem}
\begin{problem}
    Tìm max trên đường chéo phụ
\end{problem}
\begin{problem}
    Tìm max nguyên tố trong ma trận
\end{problem}
\begin{problem}
    Tìm 2 giá trị gần nhau nhất
\end{problem}
\begin{problem}
    (*) Cho ma trận vuông $A(n x n)$. Hãy tìm ma trận vuông $B(k x k)$ sao cho tổng các giá trị trên ma trận vuông này là lớn nhất
\end{problem}


\section{Kỹ thuật đếm}
\begin{problem}
    Đếm số lượng cặp giá trị đối xứng nhau qua đường chéo chính
\end{problem}
\begin{problem}
    Đếm số lượng dòng giảm
\end{problem}
\begin{problem}
    Đếm phần tử cực đại
\end{problem}
\begin{problem}
    Đếm giá trị dương trên đường chéo chính
\end{problem}
\begin{problem}
    Đếm số âm trên đường chéo phụ
\end{problem}
\begin{problem}
    Đếm số chẵn trong ma trận tam giác trên
\end{problem}


\section{Kỹ thuật đặt cờ hiệu}
\begin{problem}
    Kiểm tra đường chéo chính có tăng dần hay không
\end{problem}
\begin{problem}
    Kiểm tra ma trận có đối xứng qua chéo chính không
\end{problem}
\begin{problem}
    Kiểm tra ma trận có đối xứng qua chéo phụ không
\end{problem}
\begin{problem}
    Kiểm tra ma trận có phải là ma phương không? Ma phương là khi tổng phần tử trên các dòng, cột và 2 chéo chính phụ bằng nhau
\end{problem}


\section{Kỹ thuật sắp xếp}
\begin{problem}
    Sắp chéo chính tăng dần
\end{problem}
\begin{problem}
    Sắp chéo phụ giảm dần
\end{problem}
\begin{problem}
    Hoán vị 2 dòng
\end{problem}
\begin{problem}
    Hoán vị 2 cột
\end{problem}
\begin{problem}
    sắp các dòng tăng dần theo tổng dòng
\end{problem}
\begin{problem}
    Đưa chẵn về đầu ma trận vuông
\end{problem}
\begin{problem}
    (*) Ma trận vuông A(n x n) với n >=3. Sắp tam giác trên tăng dần từ trên xuống dưới và từ trái sang phải
\end{problem}
\begin{problem}
    (*) Ma trận vuông A với n>=3. Sắp tam giác dưới giảm dần từ trên xuống dưới và từ trái sang phải
\end{problem}
\begin{problem}
    (*) Xây dựng ma phương A
\end{problem}


\section{Các phép toán trên ma trận}
\begin{problem}
    Tổng 2 ma trận
\end{problem}
\begin{problem}
    Hiệu 2 ma trận
\end{problem}
\begin{problem}
    Tích 2 ma trận
\end{problem}
\begin{problem}
    * Ma trận nghịch đảo
\end{problem}
\begin{problem}
    * Định thức của ma trận
\end{problem}
\begin{problem}
    * Tạo ma phương bậc n x n
\end{problem}

